% ============================================================
% ODE IVP Project — LaTeX Report Template
% Compile: pdflatex report.tex  (twice for the TOC)
% Structure follows the project brief: Background -> Model ->
% Methods -> Stability -> Numerical experiments -> Conclusions.
%
% This is the MERGED template: it combines the project-specific
% guidance for Project 1 (ODE IVPs) with the course-level report
% skeleton (algorithm environment, theorem environments, MMS
% convergence table, AI Transparency Log and ICS appendices).
% ============================================================

\documentclass[11pt,a4paper]{article}

\usepackage[utf8]{inputenc}
\usepackage[T1]{fontenc}
\usepackage{amsmath, amssymb, amsthm}
\usepackage{geometry}
\geometry{a4paper, margin=2.5cm}
\usepackage{graphicx}
\usepackage{booktabs}
\usepackage{listings}
\usepackage{xcolor}
\usepackage{hyperref}
\usepackage{enumitem}
\usepackage{float}
\usepackage{caption}
\usepackage{subcaption}
\usepackage{bm}
\usepackage{algorithm}
\usepackage{algorithmic}
\usepackage{algpseudocode}
\usepackage{mdframed}

% --- colors ---
\definecolor{lecture}{RGB}{70,130,180}
\definecolor{seminar}{RGB}{205,92,92}
\definecolor{formbg}{RGB}{250,250,245}
\definecolor{formframe}{RGB}{180,170,140}

% --- code listing style (Python) ---
\lstset{
    language=Python,
    basicstyle=\ttfamily\scriptsize,
    keywordstyle=\color{blue}\bfseries,
    commentstyle=\color{green!40!black},
    stringstyle=\color{red!70!black},
    numbers=left,
    numberstyle=\tiny\color{gray},
    frame=single,
    breaklines=true,
    showstringspaces=false,
    tabsize=4,
    captionpos=b,
}

% --- theorem-like environments (course-level set) ---
\newtheorem{theorem}{Theorem}[section]
\newtheorem{lemma}[theorem]{Lemma}
\newtheorem{proposition}[theorem]{Proposition}
\newtheorem{corollary}[theorem]{Corollary}
\theoremstyle{definition}
\newtheorem{definition}[theorem]{Definition}
\newtheorem{assumption}[theorem]{Assumption}
\newtheorem{example}[theorem]{Example}
\theoremstyle{remark}
\newtheorem{remark}[theorem]{Remark}

% --- algorithm style ---
\algrenewcommand{\algorithmicrequire}{\textbf{Input:}}
\algrenewcommand{\algorithmicensure}{\textbf{Output:}}
\algrenewcommand{\algorithmiccomment}[1]{\hfill\textcolor{gray}{\# #1}}

% --- form box (for appendices) ---
\newmdenv[
    backgroundcolor=formbg,
    linecolor=formframe,
    linewidth=0.8pt,
    roundcorner=3pt,
    innermargin=5pt,
    outermargin=5pt,
    innertopmargin=8pt,
    innerbottommargin=8pt
]{formbox}

\title{\textbf{Numerical Solution of \texttt{<YOUR PROBLEM>}}\\[2mm]
\large Project 1 (ODE IVP) Report}
\author{
    Team \#\texttt{<N>} \\
    \small \texttt{<name1>, <id1> (Project Manager)} \quad
    \texttt{<name2>, <id2> (Mathematical Theory)} \\
    \small \texttt{<name3>, <id3> (Algorithm Implementation)} \quad
    \texttt{<name4>, <id4> (Visualization \& Report)} \\
    \small \texttt{<name5>, <id5> (Testing \& Validation)} \\
    \small Department of Applied Mathematics
}
\date{\today}

\begin{document}
\maketitle

\begin{abstract}
\noindent
[150--250 words answering: (1) what problem and why interesting; (2) what methods
you implemented; (3) key findings (convergence rates, stability observations);
(4) the main limitation or open question.]
\end{abstract}

\noindent\textbf{Keywords:} [e.g., Runge--Kutta methods, stiffness, von Neumann
stability, absolute stability, adaptive stepping]

\tableofcontents
\newpage

% ============================================================
\section{Introduction and Motivation}
% ============================================================

\subsection{Problem Background}
[Describe the physical, biological, or financial context of your chosen ODE
system. Why is this problem worth solving? Cite 2--3 references.]

\subsection{Mathematical Model}
[State the ODE system, initial conditions, and all parameter values. If the
system is stiff, say so and give the stiffness ratio.]

\begin{equation}
    \frac{dy}{dt} = f(t, y), \qquad y(t_0) = y_0.
    \label{eq:ivp}
\end{equation}

For the Robertson problem (example):
\begin{equation}
\begin{aligned}
    \dot{y}_1 &= -0.04\, y_1 + 10^4\, y_2 y_3, \\
    \dot{y}_2 &= \;\;0.04\, y_1 - 10^4\, y_2 y_3 - 3\times 10^7\, y_2^2, \\
    \dot{y}_3 &= \;\;3\times 10^7\, y_2^2,
\end{aligned}
\qquad y(0) = (1, 0, 0).
\label{eq:robertson}
\end{equation}

\subsection{Overview of the Report}
[A roadmap: ``In Section 2 we derive the discrete scheme and analyse stability.
In Section 3 we describe the implementation. Section 4 presents numerical
experiments. Section 5 concludes.'']

\newpage
% ============================================================
\section{Discretisation and Analysis}
% ============================================================

\subsection{Discrete Schemes}
[Derive your three methods step by step: Explicit Euler, RK4 (show the Butcher
tableau), and your implicit method with Newton. Number every equation.]
\label{sec:discrete_schemes}

We consider the autonomous initial value problem (IVP) defined on the compact temporal interval $t \in [t_0, T]$:
\begin{equation}
\label{eq:general_ivp}
\dot{\mathbf{y}}(t) = \mathbf{f}(t, \mathbf{y}(t)), \quad \mathbf{y}(t_0) = \mathbf{y}_0 \in \mathbb{R}^d,
\end{equation}
where $\mathbf{f}: [t_0, T] \times \mathbb{R}^d \to \mathbb{R}^d$ is assumed to be sufficiently smooth and locally Lipschitz continuous with respect to $\mathbf{y}$. Let the continuous interval $[t_0, T]$ be partitioned by a uniform grid $t_n = t_0 + n h$ for $n = 0, 1, \dots, N-1$, where $h = (T - t_0)/N > 0$ denotes the nominal discretization time step. Integrating \eqref{eq:general_ivp} over a single time step $[t_n, t_{n+1}]$ yields the exact Volterra integral relation:
\begin{equation}
\label{eq:integral_form}
\mathbf{y}(t_{n+1}) - \mathbf{y}(t_n) = \int_{t_n}^{t_{n+1}} \mathbf{f}(t, \mathbf{y}(t)) \, \mathrm{d}t.
\end{equation}
Different numerical approximations of the quadrature in \eqref{eq:integral_form} give rise to distinct numerical schemes with contrasting orders of accuracy, computational costs, and stability domains.

\subsubsection{Explicit Euler Method}
The simplest one-step method is obtained by approximating the integral in \eqref{eq:integral_form} via the left-hand rectangle quadrature rule, evaluating the integrand solely at the known base state $(t_n, \mathbf{y}(t_n))$:
\begin{equation}
\label{eq:explicit_euler}
\mathbf{y}_{n+1} = \mathbf{y}_n + h \mathbf{f}(t_n, \mathbf{y}_n).
\end{equation}
Equation \eqref{eq:explicit_euler} is strictly explicit; advancing the state from $t_n$ to $t_{n+1}$ incurs exactly one evaluation of the vector field $\mathbf{f}$ without requiring any linear or nonlinear system solves. While computationally inexpensive ($\mathcal{O}(d)$ floating-point operations per step), the scheme is conditionally stable and imposes severe step-size restrictions when applied to stiff ODE systems.

\subsubsection{Classical Fourth-Order Runge--Kutta Method (RK4)}
To achieve higher-order geometric fidelity without computing higher-order analytical derivatives of $\mathbf{f}$, the classical fourth-order Runge--Kutta scheme (RK4) evaluates the continuous vector field at four intermediate test stages across $[t_n, t_{n+1}]$:
\begin{align}
\mathbf{k}_1 &= \mathbf{f}(t_n, \mathbf{y}_n), \label{eq:rk4_k1} \\
\mathbf{k}_2 &= \mathbf{f}\left(t_n + \frac{h}{2}, \mathbf{y}_n + \frac{h}{2}\mathbf{k}_1\right), \label{eq:rk4_k2} \\
\mathbf{k}_3 &= \mathbf{f}\left(t_n + \frac{h}{2}, \mathbf{y}_n + \frac{h}{2}\mathbf{k}_2\right), \label{eq:rk4_k3} \\
\mathbf{k}_4 &= \mathbf{f}(t_n + h, \mathbf{y}_n + h\mathbf{k}_3), \label{eq:rk4_k4}
\end{align}
and synthesizes them via Simpson's composite quadrature weights:
\begin{equation}
\label{eq:rk4_update}
\mathbf{y}_{n+1} = \mathbf{y}_n + \frac{h}{6} \left(\mathbf{k}_1 + 2\mathbf{k}_2 + 2\mathbf{k}_3 + \mathbf{k}_4\right).
\end{equation}
The algebraic structure of any $s$-stage Runge--Kutta method can be compactly encapsulated by its Butcher tableau:
\begin{equation}
\label{eq:butcher_general}
\begin{array}{c|c}
\mathbf{c} & \mathbf{A} \\
\hline
& \mathbf{b}^T
\end{array}
\end{equation}
For classical RK4, the coefficients $\mathbf{A} \in \mathbb{R}^{4 \times 4}$, $\mathbf{b} \in \mathbb{R}^4$, and $\mathbf{c} \in \mathbb{R}^4$ are uniquely defined as:
\begin{equation}
\label{eq:rk4_butcher}
\begin{array}{c|cccc}
0   & 0   & 0   & 0 & 0 \\
1/2 & 1/2 & 0   & 0 & 0 \\
1/2 & 0   & 1/2 & 0 & 0 \\
1   & 0   & 0   & 1 & 0 \\
\hline
    & 1/6 & 1/3 & 1/3 & 1/6
\end{array}
\end{equation}
Because the Runge--Kutta matrix $\mathbf{A}$ in \eqref{eq:rk4_butcher} is strictly lower triangular with a vanishing main diagonal ($a_{ij} = 0$ for all $j \ge i$), each stage derivative $\mathbf{k}_i$ depends solely on the previously evaluated stage vectors $\mathbf{k}_1, \dots, \mathbf{k}_{i-1}$, preserving an explicit computation sequence.

\subsubsection{Implicit Euler Method with Damped Newton Iteration}
When the physical system exhibits stiffness, explicit schemes encounter numerical blow-up unless $h$ is restricted to excessively small time steps. To obtain unconditional absolute stability, we evaluate the integral in \eqref{eq:integral_form} via the right-hand rectangle rule at the future time level $(t_{n+1}, \mathbf{y}(t_{n+1}))$, yielding the Implicit (Backward) Euler scheme:
\begin{equation}
\label{eq:implicit_euler}
\mathbf{y}_{n+1} = \mathbf{y}_n + h \mathbf{f}(t_{n+1}, \mathbf{y}_{n+1}).
\end{equation}
Because the unknown state $\mathbf{y}_{n+1}$ appears nonlinearly within the argument of $\mathbf{f}$, \eqref{eq:implicit_euler} constitutes a coupled system of $d$ nonlinear algebraic equations at each time step. 

Let $\mathbf{w} \in \mathbb{R}^d$ represent the unknown state vector at $t_{n+1}$. We define the nonlinear residual mapping $\mathbf{F}: \mathbb{R}^d \to \mathbb{R}^d$ by collecting all terms on one side:
\begin{equation}
\label{eq:newton_residual}
\mathbf{F}(\mathbf{w}) := \mathbf{w} - \mathbf{y}_n - h \mathbf{f}(t_{n+1}, \mathbf{w}) = \mathbf{0}.
\end{equation}
To solve $\mathbf{F}(\mathbf{w}) = \mathbf{0}$ iteratively, we linearize $\mathbf{F}$ around a current trial estimate $\mathbf{w}^{(k)}$ via first-order multivariate Taylor expansion:
\begin{equation}
\label{eq:taylor_residual}
\mathbf{F}\left(\mathbf{w}^{(k)} + \boldsymbol{\delta}^{(k)}\right) \approx \mathbf{F}\left(\mathbf{w}^{(k)}\right) + \mathbf{J}_{\mathbf{F}}\left(\mathbf{w}^{(k)}\right) \boldsymbol{\delta}^{(k)} = \mathbf{0},
\end{equation}
where $\mathbf{J}_{\mathbf{F}}(\mathbf{w}) \in \mathbb{R}^{d \times d}$ is the Fr\'echet derivative (Jacobian matrix) of the residual operator $\mathbf{F}$ with respect to $\mathbf{w}$. Applying the multivariate chain rule to \eqref{eq:newton_residual} yields:
\begin{equation}
\label{eq:residual_jacobian}
\mathbf{J}_{\mathbf{F}}(\mathbf{w}) = \frac{\partial \mathbf{F}}{\partial \mathbf{w}} = \frac{\partial}{\partial \mathbf{w}} \left(\mathbf{w} - \mathbf{y}_n - h \mathbf{f}(t_{n+1}, \mathbf{w})\right) = \mathbf{I}_d - h \mathbf{J}_{\mathbf{f}}(t_{n+1}, \mathbf{w}),
\end{equation}
where $\mathbf{I}_d$ is the $d \times d$ identity matrix and $\mathbf{J}_{\mathbf{f}}(\cdot) = \frac{\partial \mathbf{f}}{\partial \mathbf{y}} \in \mathbb{R}^{d \times d}$ is the analytical Jacobian matrix of the continuous ODE vector field.

The standard full Newton direction $\boldsymbol{\delta}^{(k)}$ is determined by solving the linear algebraic system:
\begin{equation}
\label{eq:newton_linear_system}
\mathbf{J}_{\mathbf{F}}\left(\mathbf{w}^{(k)}\right) \boldsymbol{\delta}^{(k)} = -\mathbf{F}\left(\mathbf{w}^{(k)}\right).
\end{equation}
In practical stiff and highly nonlinear implementations, the initial state $\mathbf{w}^{(0)} = \mathbf{y}_n$ may fall outside the local quadratic attraction basin, causing standard full Newton iterations ($\mathbf{w}^{(k+1)} = \mathbf{w}^{(k)} + \boldsymbol{\delta}^{(k)}$) to suffer from severe overshoot and divergence. To enforce monotonic residual decay, we incorporate an Armijo-type backtracking line-search damping parameter $\alpha \in (0, 1]$:
\begin{equation}
\label{eq:newton_damped_update}
\mathbf{w}^{(k+1)} = \mathbf{w}^{(k)} + \alpha \boldsymbol{\delta}^{(k)}.
\end{equation}
The complete time-stepping procedure combining the outer temporal loop with the inner damped Newton solver is formalized in Algorithm~\ref{alg:implicit_euler_newton}.


\begin{algorithm}[htbp]
\caption{Implicit Euler Integration with Damped Newton Solver}
\label{alg:implicit_euler_newton}
\footnotesize % 1. 使用 footnotesize 进一步紧凑字体与行距，确保高度收缩
\begin{algorithmic}[1] % 2. 去掉末尾错误的 [H]
\STATE \textbf{Input:} Continuous RHS $\mathbf{f}(t, \mathbf{y})$, analytic Jacobian $\mathbf{J}_{\mathbf{f}}(t, \mathbf{y})$, initial state $\mathbf{y}_0$, step size $h$, final time $T$, residual tolerance $\text{tol}$, max iterations $K_{\max}$, damping shrink $\beta \in (0, 1)$.
\STATE \textbf{Output:} Discrete trajectory $\{ (t_n, \mathbf{y}_n) \}_{n=0}^N$.
\STATE $N \gets \lfloor (T - t_0) / h \rfloor$
\STATE $t \gets t_0, \quad \mathbf{y} \gets \mathbf{y}_0$
\FOR{$n = 0$ \TO $N-1$}
    \STATE $t_{\text{next}} \gets t_0 + (n + 1)h$
    \STATE $\mathbf{w} \gets \mathbf{y}_n$ \COMMENT{Initial guess: previous state}
    \STATE $\text{converged} \gets 0$
    \FOR{$k = 0$ \TO $K_{\max}-1$}
        \STATE $\mathbf{F} \gets \mathbf{w} - \mathbf{y}_n - h \mathbf{f}(t_{\text{next}}, \mathbf{w})$
        \IF{$\|\mathbf{F}\|_{\infty} < \text{tol}$}
            \STATE $\mathbf{y}_{n+1} \gets \mathbf{w}$
            \STATE $\text{converged} \gets 1$
            \STATE \textbf{break}
        \ENDIF
        
        \STATE $\mathbf{J}_{\mathbf{F}} \gets \mathbf{I}_d - h \mathbf{J}_{\mathbf{f}}(t_{\text{next}}, \mathbf{w})$
        \STATE Solve linear system $\mathbf{J}_{\mathbf{F}} \boldsymbol{\delta} = -\mathbf{F}$ for correction $\boldsymbol{\delta}$
        \STATE $\alpha \gets 1.0$ \COMMENT{Initialize full Newton step}
        \WHILE{$\|\mathbf{w} - \mathbf{y}_n - h \mathbf{f}(t_{\text{next}}, \mathbf{w} + \alpha \boldsymbol{\delta})\|_{\infty} \ge \|\mathbf{F}\|_{\infty}$ \textbf{and} $\alpha > 10^{-4}$}
            \STATE $\alpha \gets \alpha \cdot \beta$ \COMMENT{Backtrack damping factor}
        \ENDWHILE
        \STATE $\mathbf{w} \gets \mathbf{w} + \alpha \boldsymbol{\delta}$
    \ENDFOR
    \IF{$\text{converged} = 0$}
        \STATE \textbf{error} ``Newton iteration failed to converge''
    \ENDIF
\ENDFOR
\STATE \textbf{return} $\{ (t_n, \mathbf{y}_n) \}_{n=0}^N$
\end{algorithmic}
\end{algorithm}



\subsection{Consistency and Truncation Error}

\label{sec:consistency_and_truncation_error}

\subsubsection{Theoretical Assumptions and Local Defect}
To rigorously quantify the asymptotic discretization error of the numerical integrators formulated in Section~\ref{sec:discrete_schemes}, we establish the analytical foundation under classical regularity assumptions. Let the continuous vector field $\mathbf{f}: [t_0, T] \times \mathbb{R}^d \to \mathbb{R}^d$ satisfy the following conditions:
\begin{enumerate}
    \item \textbf{Lipschitz Continuity:} There exists a finite Lipschitz constant $L > 0$ such that on a bounded domain $\Omega \subset [t_0, T] \times \mathbb{R}^d$:
    \begin{equation}
    \label{eq:lipschitz_condition}
    \|\mathbf{f}(t, \mathbf{u}) - \mathbf{f}(t, \mathbf{v})\| \le L \|\mathbf{u} - \mathbf{v}\|, \quad \forall (t, \mathbf{u}), (t, \mathbf{v}) \in \Omega.
    \end{equation}
    By the Picard--Lindel\"of theorem, condition \eqref{eq:lipschitz_condition} guarantees the local existence, uniqueness, and well-posedness of the continuous trajectory $\mathbf{y}(t)$.
    \item \textbf{Sufficient Smoothness:} The theoretical solution $\mathbf{y}(t)$ is assumed to be of class $C^{p+1}[t_0, T]$, where $p \ge 1$ denotes the formal order of the numerical method, ensuring that higher-order Fr\'echet derivatives are uniformly bounded by a constant $M_{p+1} = \sup_{t \in [t_0, T]} \|\mathbf{y}^{(p+1)}(t)\| < \infty$.
\end{enumerate}

Any generic single-step integration scheme advancing from $t_n$ to $t_{n+1} = t_n + h$ can be cast into the autonomous increment form:
\begin{equation}
\label{eq:general_increment_function}
\mathbf{y}_{n+1} = \mathbf{y}_n + h \mathbf{\Phi}(t_n, \mathbf{y}_n, h),
\end{equation}
where $\mathbf{\Phi}$ represents the numerical increment function. Under the \emph{localizing assumption} (i.e., assuming the numerical solution coincides exactly with the theoretical trajectory at the current base point, $\mathbf{y}_n = \mathbf{y}(t_n)$), the one-step local defect (or local error) $\boldsymbol{\rho}_{n+1}$ is defined as:
\begin{equation}
\label{eq:local_defect_def}
\boldsymbol{\rho}_{n+1} := \mathbf{y}(t_{n+1}) - \left[ \mathbf{y}(t_n) + h \mathbf{\Phi}(t_n, \mathbf{y}(t_n), h) \right].
\end{equation}
The \emph{local truncation error} (LTE) vector per unit step, denoted by $\boldsymbol{\tau}_{n+1}$, is defined as the normalized local defect:
\begin{equation}
\label{eq:lte_def}
\boldsymbol{\tau}_{n+1} := \frac{\boldsymbol{\rho}_{n+1}}{h} = \frac{\mathbf{y}(t_{n+1}) - \mathbf{y}(t_n)}{h} - \mathbf{\Phi}(t_n, \mathbf{y}(t_n), h).
\end{equation}
A numerical scheme is defined to be \emph{consistent} of order $p$ if its local truncation error satisfies $\|\boldsymbol{\tau}_{n+1}\| = \mathcal{O}(h^p)$ as $h \to 0$, which is equivalent to requiring $\|\boldsymbol{\rho}_{n+1}\| = \mathcal{O}(h^{p+1})$.

\subsubsection{Derivation of Local Truncation Errors for Implemented Solvers}
\paragraph{1. Explicit Euler Scheme ($p=1$):}
Performing a multivariate Taylor series expansion of the exact continuous solution $\mathbf{y}(t_{n+1})$ about $t_n$ yields:
\begin{equation}
\label{eq:taylor_expansion_ee}
\mathbf{y}(t_{n+1}) = \mathbf{y}(t_n) + h \dot{\mathbf{y}}(t_n) + \frac{h^2}{2} \ddot{\mathbf{y}}(t_n) + \mathcal{O}(h^3).
\end{equation}
Recalling that $\dot{\mathbf{y}}(t_n) = \mathbf{f}(t_n, \mathbf{y}(t_n))$ and the numerical increment is $\mathbf{\Phi}_{\text{EE}} = \mathbf{f}(t_n, \mathbf{y}(t_n))$, substituting \eqref{eq:taylor_expansion_ee} into \eqref{eq:local_defect_def} leads to the exact cancellation of the zeroth- and first-order terms:
\begin{equation}
\label{eq:euler_defect_expansion}
\boldsymbol{\rho}_{n+1}^{\text{EE}} = \frac{h^2}{2} \ddot{\mathbf{y}}(t_n) + \mathcal{O}(h^3) \implies \boldsymbol{\tau}_{n+1}^{\text{EE}} = \frac{h}{2} \ddot{\mathbf{y}}(t_n) + \mathcal{O}(h^2).
\end{equation}
Hence, Explicit Euler exhibits an $\mathcal{O}(h^2)$ one-step defect and is first-order consistent ($p=1$) with leading error constant $C_2 = 1/2$.

\paragraph{2. Implicit Euler Scheme ($p=1$):}
For Implicit Euler, the increment function is $\mathbf{\Phi}_{\text{IE}} = \mathbf{f}(t_{n+1}, \mathbf{y}(t_{n+1})) = \dot{\mathbf{y}}(t_{n+1})$. Expanding $\dot{\mathbf{y}}(t_{n+1})$ backwards about $t_n$:
\begin{equation}
\label{eq:taylor_expansion_ie_slope}
\dot{\mathbf{y}}(t_{n+1}) = \dot{\mathbf{y}}(t_n) + h \ddot{\mathbf{y}}(t_n) + \mathcal{O}(h^2).
\end{equation}
Multiplying \eqref{eq:taylor_expansion_ie_slope} by $h$ and subtracting it from \eqref{eq:taylor_expansion_ee} gives the local defect:
\begin{equation}
\label{eq:ie_defect_expansion}
\boldsymbol{\rho}_{n+1}^{\text{IE}} = \mathbf{y}(t_{n+1}) - \left[ \mathbf{y}(t_n) + h \dot{\mathbf{y}}(t_{n+1}) \right] = -\frac{h^2}{2} \ddot{\mathbf{y}}(t_n) + \mathcal{O}(h^3),
\end{equation}
yielding $\boldsymbol{\tau}_{n+1}^{\text{IE}} = -\frac{h}{2} \ddot{\mathbf{y}}(t_n) + \mathcal{O}(h^2)$. The method is also first-order consistent ($p=1$) but possesses an error constant of opposite sign ($C_2 = -1/2$).

\paragraph{3. Classical Runge--Kutta Scheme (RK4, $p=4$):}
By Taylor expanding the four intermediate slope evaluations $\mathbf{k}_1, \dots, \mathbf{k}_4$ defined in \eqref{eq:rk4_k1}--\eqref{eq:rk4_k4} in terms of the elementary differentials of $\mathbf{f}$ and comparing the resultant linear combination against the continuous Taylor series of $\mathbf{y}(t_{n+1})$ up to fifth order, all algebraic order conditions up to order 4 are identically satisfied. Consequently:
\begin{equation}
\label{eq:rk4_defect}
\boldsymbol{\rho}_{n+1}^{\text{RK4}} = \boldsymbol{\psi}(t_n) h^5 + \mathcal{O}(h^6) \implies \boldsymbol{\tau}_{n+1}^{\text{RK4}} = \boldsymbol{\psi}(t_n) h^4 + \mathcal{O}(h^5),
\end{equation}
where $\boldsymbol{\psi}(t_n)$ denotes the principal error function composed of elementary differentials of order 5. Hence, RK4 is consistent of order $p=4$.

\subsubsection{Adaptive Step-Size Control via Step-Doubling Richardson Extrapolation}
Practical simulation of multiscale or stiff dynamical systems necessitates dynamic step-size selection to maintain uniform accuracy while minimizing computational overhead. To estimate the local truncation error without relying on analytical higher derivatives, we employ the \emph{step-doubling} (Richardson extrapolation) strategy.

Let the base numerical integrator have order $p$. From a shared accepted state $(t_n, \mathbf{y}_n)$, we compute two independent numerical predictions at $t_{n+1} = t_n + h$:
\begin{enumerate}
    \item \textbf{Coarse Solution ($\mathbf{y}_c$):} A single trial step spanning the full increment $h$:
    \begin{equation}
    \label{eq:coarse_step}
    \mathbf{y}_c = \mathbf{y}(t_n + h) - \boldsymbol{\psi}(t_n) h^{p+1} + \mathcal{O}(h^{p+2}).
    \end{equation}
    \item \textbf{Fine Solution ($\mathbf{y}_f$):} Two consecutive sub-steps spanning $h/2$. Assuming the principal error coefficient $\boldsymbol{\psi}(t)$ remains locally constant across $[t_n, t_n + h]$, the defects from both half-steps superimpose linearly:
    \begin{align}
    \mathbf{y}_{1/2} &= \mathbf{y}\left(t_n + \frac{h}{2}\right) - \boldsymbol{\psi}(t_n) \left(\frac{h}{2}\right)^{p+1} + \mathcal{O}(h^{p+2}), \label{eq:fine_half_step} \\
    \mathbf{y}_f &= \mathbf{y}(t_n + h) - 2 \boldsymbol{\psi}(t_n) \left(\frac{h}{2}\right)^{p+1} + \mathcal{O}(h^{p+2}) = \mathbf{y}(t_n + h) - \frac{\boldsymbol{\psi}(t_n) h^{p+1}}{2^p} + \mathcal{O}(h^{p+2}). \label{eq:fine_step}
    \end{align}
\end{enumerate}
Subtracting \eqref{eq:fine_step} from \eqref{eq:coarse_step} eliminates the unknown exact continuous state $\mathbf{y}(t_n + h)$:
\begin{equation}
\label{eq:richardson_subtraction}
\mathbf{y}_c - \mathbf{y}_f = -\boldsymbol{\psi}(t_n) h^{p+1} \left(1 - \frac{1}{2^p}\right) + \mathcal{O}(h^{p+2}) = -\frac{2^p - 1}{2^p} \boldsymbol{\psi}(t_n) h^{p+1} + \mathcal{O}(h^{p+2}).
\end{equation}
Solving \eqref{eq:richardson_subtraction} directly isolates the leading error term:
\begin{equation}
\label{eq:principal_error_isolated}
\frac{\boldsymbol{\psi}(t_n) h^{p+1}}{2^p} = \frac{\mathbf{y}_f - \mathbf{y}_c}{2^p - 1} + \mathcal{O}(h^{p+2}).
\end{equation}
Consequently, the computable local error estimate of the more accurate fine solution $\mathbf{y}_f$ is given by:
\begin{equation}
\label{eq:local_error_estimator}
\hat{\mathbf{e}}_n := \mathbf{y}(t_n + h) - \mathbf{y}_f \approx \frac{\mathbf{y}_f - \mathbf{y}_c}{2^p - 1}.
\end{equation}
For our implemented solvers, the denominator evaluates explicitly to:
\begin{itemize}
    \item \textbf{Implicit / Explicit Euler ($p=1$):} $2^1 - 1 = 1 \implies \hat{\mathbf{e}}_n \approx \mathbf{y}_f - \mathbf{y}_c$.
    \item \textbf{Classical RK4 ($p=4$):} $2^4 - 1 = 15 \implies \hat{\mathbf{e}}_n \approx \frac{\mathbf{y}_f - \mathbf{y}_c}{15}$.
\end{itemize}

\subsubsection{Error Normalization and Step-Update Controller}
To render the acceptance decision scale-invariant across state variables of disparate magnitudes, we define a component-wise tolerance vector $\mathbf{s} \in \mathbb{R}^d$ combining user-prescribed absolute ($\text{atol}$) and relative ($\text{rtol}$) tolerances:
\begin{equation}
\label{eq:scale_vector}
s_i = \text{atol}_i + \text{rtol} \cdot \max\left(|y_{n, i}|, \, |y_{f, i}|\right), \quad i = 1, \dots, d.
\end{equation}
The overall scalar normalized error ratio $E$ is evaluated using the weighted $L_\infty$ norm:
\begin{equation}
\label{eq:normalized_error_ratio}
E := \max_{1 \le i \le d} \left| \frac{\hat{e}_{n, i}}{s_i} \right|.
\end{equation}
The operational decision logic is partitioned into two mutually exclusive regimes:
\begin{itemize}
    \item \textbf{Step Acceptance ($E \le 1$):} The estimated error falls strictly within the permissible tolerance corridor. The time level and state are advanced using \emph{local extrapolation}:
    \begin{equation}
    \label{eq:local_extrapolation}
    t_{n+1} \gets t_n + h, \quad \mathbf{y}_{n+1} \gets \mathbf{y}_f.
    \end{equation}
    \item \textbf{Step Rejection ($E > 1$):} The local error exceeds the target envelope. In accordance with strict conservation, \emph{a rejected step has not occurred}; neither $t_n$ nor $\mathbf{y}_n$ is modified, and the step is immediately recomputed from $(t_n, \mathbf{y}_n)$ using an attenuated step size.
\end{itemize}

Under the asymptotic error relation $E(h) \propto h^{p+1}$, setting the target normalized error for the prospective step to unity ($E(h_{\text{opt}}) \approx 1$) yields the unconstrained scale factor $r = E^{-1/(p+1)}$. To prevent step-size hunting and severe limit-cycle oscillations induced by higher-order nonlinear transients, we incorporate a conservative safety margin $S \in [0.8, 0.9]$ alongside hard growth and shrinkage limiter bounds $[r_{\min}, r_{\max}]$:
\begin{equation}
\label{eq:step_update_formula}
h_{\text{new}} = h \cdot \min\left(r_{\max}, \, \max\left(r_{\min}, \, S \cdot E^{-\frac{1}{p+1}}\right)\right),
\end{equation}
where typical robust parameters are set to $S = 0.9$, $r_{\min} = 0.2$, and $r_{\max} = 2.0$.

\subsection{Stability Analysis}

\begin{equation}
    R_{\text{Euler}}(z) = 1 + z, \qquad
    R_{\text{RK4}}(z) = 1 + z + \frac{z^2}{2} + \frac{z^3}{6} + \frac{z^4}{24},
    \qquad
    R_{\text{impl.Euler}}(z) = \frac{1}{1 - z}.
\end{equation}

\label{sec:stability_analysis}

While consistency governs the asymptotic error behavior as $h \to 0$, the numerical feasibility of an integrator at finite step sizes is governed by absolute stability. Stiff dynamical regimes impose severe step-size restrictions on explicit integrators, rendering stability, rather than accuracy, the primary bottleneck.

\subsubsection{Scalar Test Problem and Amplification Factors}
We consider the classical Dahlquist scalar linear test equation:
\begin{equation}
\label{eq:dahlquist_test}
\dot{y}(t) = \lambda y(t), \quad y(0) = y_0, \quad \lambda \in \mathbb{C}, \quad \mathrm{Re}(\lambda) < 0,
\end{equation}
whose exact continuous trajectory decays monotonically: $|y(t)| = |y_0| e^{\mathrm{Re}(\lambda)t} \to 0$ as $t \to \infty$. Introducing the dimensionless complex parameter $z := h\lambda \in \mathbb{C}$, applying a generic single-step numerical integrator to \eqref{eq:dahlquist_test} yields the discrete geometric recurrence:
\begin{equation}
\label{eq:discrete_recurrence}
y_{n+1} = R(z) y_n = \left[R(z)\right]^{n+1} y_0,
\end{equation}
where $R(z): \mathbb{C} \to \mathbb{C}$ denotes the rational or polynomial stability amplification function. The discrete perturbation sequence $\{y_n\}_{n=0}^\infty$ remains uniformly bounded if and only if the complex amplification factor satisfies the von Neumann criterion:
\begin{equation}
\label{eq:stability_criterion}
|R(z)| \le 1.
\end{equation}
The region of absolute stability in the complex plane is formally defined as:
\begin{equation}
\label{eq:stability_region_def}
\mathcal{S} := \left\{ z \in \mathbb{C} : |R(z)| \le 1 \right\}.
\end{equation}

We derive the analytic form of $R(z)$ for the three implemented solvers:
\begin{enumerate}
    \item \textbf{Explicit Euler:} Substituting $f(y) = \lambda y$ into \eqref{eq:explicit_euler} gives:
    \begin{equation}
    \label{eq:R_EE}
    y_{n+1} = y_n + h\lambda y_n = (1 + z) y_n \implies R_{\mathrm{EE}}(z) = 1 + z.
    \end{equation}
    Condition \eqref{eq:stability_criterion} requires $|1 + z| \le 1$. In the complex plane $z = x + \mathrm{i}y$, this delineates a closed circular disk of unit radius centered at $(-1, 0)$:
    \begin{equation}
    \label{eq:EE_boundary}
    (x + 1)^2 + y^2 \le 1.
    \end{equation}
    For purely negative real eigenvalues ($\lambda = -a$ with $a > 0$, $y = 0$), the stability interval is strictly restricted to $x \in [-2, 0]$, dictating the classical step-size ceiling:
    \begin{equation}
    \label{eq:h_crit_EE}
    h \le \frac{2}{|\lambda|}.
    \end{equation}
    Crucially, along the imaginary axis ($x = 0, y \ne 0$), $|R_{\mathrm{EE}}(\mathrm{i}y)|^2 = 1 + y^2 > 1$ for all $y \ne 0$. Hence, Explicit Euler contains no segment of the imaginary axis and is unconditionally unstable for purely oscillatory systems.

    \item \textbf{Classical Fourth-Order Runge--Kutta (RK4):} Applying RK4 stage equations \eqref{eq:rk4_k1}--\eqref{eq:rk4_k4} recursively to $\dot{y} = \lambda y$ yields the stage slopes:
    \begin{align}
    h k_1 &= z y_n, \label{eq:rk4_stg1} \\
    h k_2 &= z \left( y_n + \frac{1}{2} h k_1 \right) = \left( z + \frac{z^2}{2} \right) y_n, \label{eq:rk4_stg2} \\
    h k_3 &= z \left( y_n + \frac{1}{2} h k_2 \right) = \left( z + \frac{z^2}{2} + \frac{z^3}{4} \right) y_n, \label{eq:rk4_stg3} \\
    h k_4 &= z \left( y_n + h k_3 \right) = \left( z + z^2 + \frac{z^3}{2} + \frac{z^4}{4} \right) y_n. \label{eq:rk4_stg4}
    \end{align}
    Synthesizing into the final update \eqref{eq:rk4_update} leads directly to the truncated Taylor polynomial of the continuous exponential map $e^z$:
    \begin{equation}
    \label{eq:R_RK4}
    R_{\mathrm{RK4}}(z) = 1 + z + \frac{z^2}{2!} + \frac{z^3}{3!} + \frac{z^4}{4!}.
    \end{equation}
    The stability boundary $|R_{\mathrm{RK4}}(z)| = 1$ encloses an extended domain shown in Figure~\ref{fig:stability_overlay}. Along the negative real axis ($z = -a, a > 0$), setting $R_{\mathrm{RK4}}(-a) = 1$ yields:
    \begin{equation}
    \label{eq:rk4_real_root}
    -a + \frac{a^2}{2} - \frac{a^3}{6} + \frac{a^4}{24} = 0 \implies a\left(a^3 - 4a^2 + 12a - 24\right) = 0.
    \end{equation}
    Factoring out the trivial root $a = 0$, the cubic factor has a unique real root at $a^* \approx 2.78529$, establishing the negative real stability interval $\mathcal{S}_{\mathrm{RK4}} \cap \mathbb{R} = [-2.78529, 0]$. Along the imaginary axis $z = \mathrm{i}y$, $|R_{\mathrm{RK4}}(\mathrm{i}y)|^2 \le 1$ holds for $|y| \le 2\sqrt{2} \approx 2.8284$, conferring robust stability for moderately oscillatory modes.

    \item \textbf{Implicit Euler:} Evaluating at the implicit state $y_{n+1}$ yields:
    \begin{equation}
    \label{eq:R_IE}
    y_{n+1} = y_n + z y_{n+1} \implies (1 - z)y_{n+1} = y_n \implies R_{\mathrm{IE}}(z) = \frac{1}{1 - z}.
    \end{equation}
    The stability criterion $|R_{\mathrm{IE}}(z)| \le 1$ is equivalent to $|1 - z| \ge 1$, which defines the exterior of the open unit disk centered at $(+1, 0)$ in the complex plane:
    \begin{equation}
    \label{eq:IE_domain}
    (x - 1)^2 + y^2 \ge 1.
    \end{equation}
\end{enumerate}

\subsubsection{A-Stability and L-Stability Properties}
A numerical integration scheme is defined to be \emph{A-stable} if its region of absolute stability contains the entire open left complex half-plane:
\begin{equation}
\label{eq:a_stability_def}
\mathbb{C}^- := \{ z \in \mathbb{C} : \mathrm{Re}(z) \le 0 \} \subseteq \mathcal{S}.
\end{equation}
For Implicit Euler, let $z = x + \mathrm{i}y$ with $x \le 0$. The squared modulus evaluates to:
\begin{equation}
\label{eq:a_stability_proof}
|1 - z|^2 = (1 - x)^2 + y^2 = 1 - 2x + x^2 + y^2.
\end{equation}
Since $x \le 0$, $-2x \ge 0$, and $x^2 + y^2 \ge 0$, it follows that $|1 - z|^2 \ge 1$ identically for all $z \in \mathbb{C}^-$. Therefore, Implicit Euler is unconditionally A-stable.

Furthermore, an A-stable method is classified as \emph{L-stable} if its stability function satisfies the asymptotic damping limit at stiff infinity:
\begin{equation}
\label{eq:l_stability_def}
\lim_{|z| \to \infty} |R(z)| = 0.
\end{equation}
Evaluating \eqref{eq:R_IE} gives $\lim_{|z| \to \infty} \left|\frac{1}{1 - z}\right| = 0$, confirming that Implicit Euler is L-stable. Physically, this property guarantees that infinitely stiff transient decay modes ($\mathrm{Re}(\lambda) \ll 0$) are annihilated in a single time step without inducing parasitic high-frequency numerical oscillations. Conversely, explicit polynomials $|R_{\mathrm{EE}}(z)| \to \infty$ and $|R_{\mathrm{RK4}}(z)| \to \infty$ as $|z| \to \infty$; by Dahlquist's barrier theorems, explicit Runge--Kutta schemes cannot be A-stable.

\subsubsection{Problem-Specific Linearisation and Eigenvalue Spectrum (Topic 5)}
To evaluate practical numerical stability beyond generic scalar benchmarks, we specialize the analysis to Topic~5: the polar-factor gradient flow on $\mathbb{R}^{n \times n}$:
\begin{equation}
\label{eq:polar_flow_spec}
\dot{\mathbf{X}} = \mathbf{f}(\mathbf{X}) = \mathbf{X}\left(\mathbf{I}_n - \mathbf{X}^T \mathbf{X}\right), \quad \mathbf{X}(0) = \mathbf{X}_0.
\end{equation}
Let $\mathbf{X}(t) = \mathbf{U} \boldsymbol{\Sigma}(t) \mathbf{V}^T$ denote the singular value decomposition (SVD), where $\boldsymbol{\Sigma}(t) = \mathrm{diag}(s_1(t), \dots, s_n(t))$ with $s_i(t) > 0$. Because the singular vectors $\mathbf{U}, \mathbf{V} \in \mathcal{O}(n)$ remain invariant under the flow, \eqref{eq:polar_flow_spec} decouples into $n$ independent scalar singular value differential equations:
\begin{equation}
\label{eq:singular_value_ode}
\dot{s}_i(t) = s_i(t) \left(1 - s_i^2(t)\right), \quad i = 1, \dots, n.
\end{equation}
Linearizing \eqref{eq:singular_value_ode} around any instantaneous singular state $s_i$ yields the local scalar Jacobian:
\begin{equation}
\label{eq:singular_jacobian}
\lambda_i(s_i) = \frac{\partial \dot{s}_i}{\partial s_i} = 1 - 3 s_i^2.
\end{equation}
Near the asymptotic orthogonal equilibrium $\mathbf{X}^* = \mathbf{U}\mathbf{V}^T$ ($s_i^* = 1$), the linearized eigenvalue is $\lambda_i(1) = -2$, confirming asymptotic exponential contraction. However, during the initial transient phase where the matrix contains amplified singular values $s_{\max} > 1$, the continuous mode decays with severe stiffness:
\begin{equation}
\label{eq:lambda_min_transient}
\lambda_{\mathrm{stiff}} = 1 - 3 s_{\max}^2 \ll 0.
\end{equation}
For our benchmark matrix possessing an initial maximum singular value $s_{\max}(0) = 3.0$, the instantaneous stiffness eigenvalue reaches $\lambda_{\mathrm{fast}} = 1 - 3(3.0)^2 = -26.0$. Meanwhile, modes with $s_{\min} \approx 1$ possess eigenvalues near $-2.0$, yielding a localized transient stiffness ratio:
\begin{equation}
\label{eq:stiffness_ratio_topic5}
\mathrm{SR}(t_0) = \frac{\max_i |\mathrm{Re}(\lambda_i)|}{\min_i |\mathrm{Re}(\lambda_i)|} = \frac{|-26.0|}{|-2.0|} = 13.0.
\end{equation}

For the full vectorized matrix field $\mathbf{y} = \mathrm{vec}(\mathbf{X}) \in \mathbb{R}^{n^2}$, the Fr\'echet derivative in matrix direction $\mathbf{H} \in \mathbb{R}^{n \times n}$ is given by:
\begin{equation}
\label{eq:frechet_derivative}
\mathcal{D}\mathbf{f}(\mathbf{X})[\mathbf{H}] = \mathbf{H}\left(\mathbf{I}_n - \mathbf{X}^T \mathbf{X}\right) - \mathbf{X}\left(\mathbf{H}^T \mathbf{X} + \mathbf{X}^T \mathbf{H}\right).
\end{equation}
Vectorizing \eqref{eq:frechet_derivative} using the Kronecker identity $\mathrm{vec}(\mathbf{A}\mathbf{B}\mathbf{C}) = (\mathbf{C}^T \otimes \mathbf{A})\mathrm{vec}(\mathbf{B})$ and the commutation matrix $\mathbf{K}_{n,n}\mathrm{vec}(\mathbf{H}) = \mathrm{vec}(\mathbf{H}^T)$ yields the exact $n^2 \times n^2$ analytical Jacobian:
\begin{equation}
\label{eq:analytic_jacobian_kron}
\mathbf{J}_{\mathbf{f}}(\mathbf{X}) = \mathbf{I}_n \otimes \left(\mathbf{I}_n - \mathbf{X}\mathbf{X}^T\right) - \left(\mathbf{X}^T \otimes \mathbf{X}\right)\mathbf{K}_{n,n} - \left(\mathbf{X}^T \mathbf{X}\right)^T \otimes \mathbf{I}_n.
\end{equation}
Evaluating the spectrum $\sigma(\mathbf{J}_{\mathbf{f}}(\mathbf{X}_0))$ reveals that all physical eigenvalues lie strictly on the negative real line $\lambda_j \in [-26.0, -2.0]$, confirming a dissipative spectrum free of imaginary components.

\subsubsection{Theoretical Step-Size Limits vs. Empirical Disclaimers}
Mapping the extreme spectral point $\lambda_{\min} = -26.0$ onto the real stability intervals of the explicit solvers yields rigorous theoretical upper bounds for non-divergent integration:
\begin{align}
\text{Explicit Euler:} \quad & h_{\mathrm{crit}}^{\mathrm{EE}} = \frac{2}{|\lambda_{\min}|} = \frac{2}{26.0} \approx 0.0769, \label{eq:hcrit_EE_calc} \\
\text{Classical RK4:} \quad & h_{\mathrm{crit}}^{\mathrm{RK4}} = \frac{2.78529}{|\lambda_{\min}|} = \frac{2.78529}{26.0} \approx 0.1071. \label{eq:hcrit_RK4_calc}
\end{align}
Figure~\ref{fig:stability_overlay} illustrates the absolute stability domains of the three methods overlaid with the scaled eigenvalues $z_j = h\lambda_j$ at various representative step sizes.

\paragraph{Mandatory Analytical Disclaimer:}
\emph{We formally emphasize that the $h\lambda$ spectral overlay represents a local frozen-Jacobian diagnostic rather than a global nonlinear stability proof. Because the Jacobian $\mathbf{J}_{\mathbf{f}}(\mathbf{X}(t))$ varies continuously along the matrix trajectory, linearized eigenvalue containment $|R(h\lambda_j)| \le 1$ is a necessary local condition but insufficient to guarantee nonlinear convergence. Severe non-normal matrix transient interactions or finite-amplitude numerical perturbations can trigger catastrophic nonlinear instability even when frozen eigenvalues lie nominally within $\mathcal{S}$. Consequently, theoretical stability boundaries \eqref{eq:hcrit_EE_calc}--\eqref{eq:hcrit_RK4_calc} must be empirically verified via systematic step-size sweeps and finite perturbation experiments in Section~\ref{sec:numerical_experiments}.}

\subsection{Convergence}

\label{sec:convergence}

Having established the consistency orders in Section~\ref{sec:consistency_and_truncation_error} and linear stability boundaries in Section~\ref{sec:stability_analysis}, we now unify these components to prove the global convergence of the numerical trajectory towards the theoretical solution.

\subsubsection{Discrete Error Recurrence and the Discrete Gr\"onwall Lemma}
Let $\mathbf{y}(t) \in C^{p+1}[t_0, T]$ be the exact continuous trajectory of the initial value problem \eqref{eq:general_ivp}, and let $\{\mathbf{y}_n\}_{n=0}^N$ be the discrete numerical approximation generated by the one-step scheme:
\begin{equation}
\label{eq:increment_form_convergence}
\mathbf{y}_{n+1} = \mathbf{y}_n + h \mathbf{\Phi}(t_n, \mathbf{y}_n, h), \quad \mathbf{y}_0 = \mathbf{y}(t_0),
\end{equation}
where $h = (T - t_0)/N$. The true solution identically satisfies the perturbed relation:
\begin{equation}
\label{eq:exact_perturbed_relation}
\mathbf{y}(t_{n+1}) = \mathbf{y}(t_n) + h \mathbf{\Phi}(t_n, \mathbf{y}(t_n), h) + h \boldsymbol{\tau}_{n+1},
\end{equation}
where $\boldsymbol{\tau}_{n+1}$ is the local truncation error vector defined in \eqref{eq:lte_def}. We assume the numerical increment function $\mathbf{\Phi}(t, \mathbf{u}, h)$ inherits the Lipschitz continuity of the continuous vector field $\mathbf{f}$ with an effective Lipschitz constant $\Lambda > 0$:
\begin{equation}
\label{eq:increment_lipschitz}
\|\mathbf{\Phi}(t, \mathbf{u}, h) - \mathbf{\Phi}(t, \mathbf{v}, h)\| \le \Lambda \|\mathbf{u} - \mathbf{v}\|, \quad \forall (t, \mathbf{u}), (t, \mathbf{v}) \in [t_0, T] \times \mathbb{R}^d, \; \forall h \in (0, h_0].
\end{equation}
We define the global discretization error at grid node $t_n$ as:
\begin{equation}
\label{eq:global_error_def}
\mathbf{e}_n := \mathbf{y}(t_n) - \mathbf{y}_n, \quad n = 0, 1, \dots, N.
\end{equation}
Subtracting \eqref{eq:increment_form_convergence} from \eqref{eq:exact_perturbed_relation} yields the dynamic recurrence for the error vector:
\begin{equation}
\label{eq:error_recurrence_vector}
\mathbf{e}_{n+1} = \mathbf{e}_n + h \left[ \mathbf{\Phi}(t_n, \mathbf{y}(t_n), h) - \mathbf{\Phi}(t_n, \mathbf{y}_n, h) \right] + h \boldsymbol{\tau}_{n+1}.
\end{equation}
Taking Euclidean norms on both sides of \eqref{eq:error_recurrence_vector} and applying the triangle inequality alongside \eqref{eq:increment_lipschitz}:
\begin{align}
\|\mathbf{e}_{n+1}\| &\le \|\mathbf{e}_n\| + h \|\mathbf{\Phi}(t_n, \mathbf{y}(t_n), h) - \mathbf{\Phi}(t_n, \mathbf{y}_n, h)\| + h \|\boldsymbol{\tau}_{n+1}\| \nonumber \\
&\le \|\mathbf{e}_n\| + h \Lambda \|\mathbf{y}(t_n) - \mathbf{y}_n\| + h \|\boldsymbol{\tau}_{n+1}\| \nonumber \\
&= (1 + h \Lambda) \|\mathbf{e}_n\| + h \|\boldsymbol{\tau}_{n+1}\|. \label{eq:scalar_error_inequality}
\end{align}
Let $\tau_{\max}(h) := \max_{1 \le k \le N} \|\boldsymbol{\tau}_k\|$. Unrolling the scalar recurrence \eqref{eq:scalar_error_inequality} inductively from the initial state yields:
\begin{equation}
\label{eq:unrolled_error_sum}
\|\mathbf{e}_n\| \le (1 + h\Lambda)^n \|\mathbf{e}_0\| + h \tau_{\max}(h) \sum_{j=0}^{n-1} (1 + h\Lambda)^j.
\end{equation}
Assuming exact initial data ($\|\mathbf{e}_0\| = 0$) and summing the finite geometric progression with ratio $1 + h\Lambda > 1$:
\begin{equation}
\label{eq:geometric_sum_evaluation}
\sum_{j=0}^{n-1} (1 + h\Lambda)^j = \frac{(1 + h\Lambda)^n - 1}{(1 + h\Lambda) - 1} = \frac{(1 + h\Lambda)^n - 1}{h\Lambda}.
\end{equation}
Substituting \eqref{eq:geometric_sum_evaluation} into \eqref{eq:unrolled_error_sum}, and noting the fundamental inequality $1 + x \le e^x$ for all $x \ge 0$, we have:
\begin{equation}
\label{eq:exponential_upper_bound}
(1 + h\Lambda)^n \le \left(e^{h\Lambda}\right)^n = e^{n h \Lambda} = e^{\Lambda (t_n - t_0)} \le e^{\Lambda(T - t_0)}.
\end{equation}
Combining \eqref{eq:unrolled_error_sum}--\eqref{eq:exponential_upper_bound} produces the classical continuous Gr\"onwall bound on the global discretization error:
\begin{equation}
\label{eq:global_convergence_bound}
\|\mathbf{e}_n\| \le \frac{e^{\Lambda(T - t_0)} - 1}{\Lambda} \tau_{\max}(h), \quad \forall n = 0, 1, \dots, N.
\end{equation}

\subsubsection{Global Convergence Order and the Lax--Dahlquist Equivalence}
Theorem~\eqref{eq:global_convergence_bound} provides the fundamental mathematical bridge between the local truncation error rate $\boldsymbol{\tau}$ and the global error $\mathbf{e}_n$:
\begin{itemize}
    \item \textbf{Explicit / Implicit Euler ($p=1$):} Since $\|\boldsymbol{\tau}_{n+1}\| \le C_1 h$, the terminal error satisfies:
    \begin{equation}
    \label{eq:euler_global_order}
    \|\mathbf{e}_N\| \le \left( \frac{e^{\Lambda(T - t_0)} - 1}{\Lambda} C_1 \right) h = \mathcal{O}(h).
    \end{equation}
    \item \textbf{Classical Runge--Kutta (RK4, $p=4$):} Since $\|\boldsymbol{\tau}_{n+1}\| \le C_4 h^4$, the terminal error satisfies:
    \begin{equation}
    \label{eq:rk4_global_order}
    \|\mathbf{e}_N\| \le \left( \frac{e^{\Lambda(T - t_0)} - 1}{\Lambda} C_4 \right) h^4 = \mathcal{O}(h^4).
    \end{equation}
\end{itemize}
This explains the apparent paradox regarding order indices: while an individual integration step generates a local defect of order $\mathcal{O}(h^{p+1})$, the accumulation of $N = \mathcal{O}(h^{-1})$ steps over a compact temporal interval $[t_0, T]$ causes the discrete sequence to converge globally at order $\mathcal{O}(h^p)$.

In modern numerical analysis, this result is generalized by the celebrated \emph{Lax--Dahlquist Equivalence Theorem} adapted for one-step methods:
\begin{equation}
\label{eq:lax_dahlquist_relation}
\text{\textbf{Consistency}} \; (\boldsymbol{\tau} \to \mathbf{0}) \quad + \quad \text{\textbf{Zero-Stability}} \quad \Longleftrightarrow \quad \text{\textbf{Convergence}} \; (\|\mathbf{e}_N\| \to 0).
\end{equation}
A numerical scheme is \emph{zero-stable} if small perturbations in the initial conditions or intermediate arithmetic operations remain uniformly bounded as $h \to 0$ over a finite interval. For any generic one-step integrator \eqref{eq:increment_form_convergence}, the characteristic polynomial of the difference scheme in the limit $h \to 0$ degenerates strictly to:
\begin{equation}
\label{eq:root_condition_onestep}
\rho(\zeta) = \zeta - 1 = 0.
\end{equation}
The root $\zeta_1 = 1$ is simple and lies on the complex unit circle ($|\zeta_1| \le 1$), identically satisfying the Dahlquist root condition. Consequently, \emph{all consistent one-step integration schemes are unconditionally zero-stable}. Global convergence is thus guaranteed entirely by the algebraic consistency order of the numerical scheme, provided the continuous vector field $\mathbf{f}$ remains locally Lipschitz continuous.

\subsubsection{Implications for Empirical Log-Log Validation}
Taking the natural logarithm of the asymptotic global error relation $\|\mathbf{e}_N(h)\| \approx C h^p$ yields:
\begin{equation}
\label{eq:loglog_linear_relation}
\ln \|\mathbf{e}_N(h)\| \approx p \ln h + \ln C.
\end{equation}
Equation \eqref{eq:loglog_linear_relation} establishes that plotting the terminal error against the discretization step size on a logarithmic scale yields an affine straight line whose asymptotic slope is identical to the theoretical convergence order $p$. In Section~\ref{sec:numerical_experiments}, this linear relationship serves as our primary quantitative metric to corroborate the theoretical convergence orders of Explicit Euler ($p \approx 1$), Implicit Euler ($p \approx 1$), and Classical RK4 ($p \approx 4$).

\newpage
% ============================================================
\section{Implementation}
% ============================================================

\subsection{Algorithm}
[Present the core algorithm in pseudocode using the \texttt{algorithmic}
environment.]

\begin{algorithm}[H]
\caption{[Your method name]}
\label{alg:main}
\begin{algorithmic}[1]
\Require [initial data $y_0$, step size $h$, final time $T$]
\Ensure [approximate solution at $T$]
\State [Step 1]
\For{$n = 0, 1, \ldots, N-1$}
    \State [Loop body]
    \If{[condition]}
        \State [Action]
    \EndIf
\EndFor
\State \Return [result]
\end{algorithmic}
\end{algorithm}

\subsection{Code Structure}
[Briefly describe the organisation of your codebase. A table of files and their
purposes is helpful.]

\begin{table}[H]
\centering
\begin{tabular}{ll}
\toprule
\textbf{File} & \textbf{Purpose} \\
\midrule
\texttt{methods.py} & Core time-stepping: Euler, RK4, implicit Euler with Newton \\
\texttt{model.py} & The right-hand side $f(t,y)$ and its Jacobian \\
\texttt{adaptive.py} & Adaptive step-size controller \\
\texttt{run\_all.py} & Reproduces all report figures \\
\texttt{validation\_tests.py} & PASS/FAIL test suite \\
\bottomrule
\end{tabular}
\end{table}

\subsection{Key Design Decisions}
[Discuss non-obvious choices: why this implicit method, how Newton's initial
guess is chosen, how the adaptive controller estimates local error.]

\newpage
% ============================================================
\section{Numerical Experiments}
% ============================================================

\subsection{Verification: Method of Manufactured Solutions or Exact Solutions}
[Present a test case with a known solution. Show the error table and the
convergence plot.]

\begin{table}[H]
\centering
\begin{tabular}{cccc}
\toprule
\textbf{Step size $h$} & \textbf{Error} & \textbf{Ratio} & \textbf{Order} \\
\midrule
$0.1$    & [value] & ---  & ---  \\
$0.05$   & [value] & [ratio] & [order] \\
$0.025$  & [value] & [ratio] & [order] \\
\bottomrule
\end{tabular}
\caption{Convergence study for [test case].}
\end{table}

\subsection{Stability Experiments}
[Show that your explicit methods blow up beyond their stability limit while the
implicit method does not. Compare with the predicted threshold.]

\subsection{Adaptivity}
[Show the step size actually changes, the local error stays near the tolerance,
and the $h(t)$ trajectory.]

\subsection{Main Results and Comparison of Methods}
[Present the experiments that answer the project's central question. Use
high-quality figures with interpretive captions.]

\begin{figure}[H]
\centering
\fbox{\parbox{0.8\textwidth}{\centering \vspace{3em} [Insert your figure here] \vspace{3em}}}
\caption{[Caption that interprets, not merely describes. E.g., ``The log-log plot
confirms second-order convergence (slope 1.98) on the smooth test case, and shows
the round-off floor at $h \approx 10^{-6}$.'']}
\label{fig:convergence}
\end{figure}

\subsection{Limitations and Robustness}
[Discuss parameter regimes where your method struggles. This is not a weakness ---
it is scientific honesty.]

\newpage
% ============================================================
\section{Conclusions}
% ============================================================

[Summarise in 3--4 bullet points:]
\begin{itemize}[leftmargin=2em]
\item [What you achieved.]
\item [The most surprising or instructive finding.]
\item [The main limitation you identified.]
\item [A concrete direction for future work.]
\end{itemize}

\newpage
% ============================================================
\section*{References}
% ============================================================

[Use BibTeX or list references manually. Include at least: (1) the textbook from
which you learned the method; (2) a paper applying the method to a problem
similar to yours; (3) a software reference (NumPy, SciPy, scikit-fem).]

\begin{thebibliography}{9}
\bibitem{ref1} Author, A. \emph{Title}. Journal, year.
\bibitem{ref2} Author, B. \emph{Book title}. Publisher, year.
\end{thebibliography}

\newpage
% ============================================================
\appendix
\section{AI Transparency Log}
% ============================================================

[Append the completed team AI Transparency Log here. See
\texttt{06\_Templates/ICS} for the full form: tools used, significant claim audits,
what you changed and why, and the intellectual ownership declarations.]

\section{Individual Contribution Statements}
% ============================================================

Each team member must complete the following (repeat for all 5 members):

\begin{formbox}
\textbf{Name:} \underline{\hspace{5cm}}

\vspace{0.5em}
\textbf{1. My specific contributions to this project:}

[3--5 sentences: which equations you derived, which code modules you wrote,
which figures you produced, which report sections you drafted.]

\vspace{0.5em}
\textbf{2. The hardest conceptual obstacle I encountered:}

[One specific mathematical or algorithmic difficulty and how you overcame it.]

\vspace{0.5em}
\textbf{3. One piece of AI-generated output my team used, and what I changed:}

[Cross-reference an item from the AI Transparency Log.]

\vspace{0.5em}
\textbf{4. One thing I still do not fully understand:}

[Not penalised; helps the instructor plan the next project.]

\vspace{0.5em}
\textbf{Signature:} \underline{\hspace{4cm}} \quad \textbf{Date:} \underline{\hspace{2cm}}
\end{formbox}

\section{Supplementary Figures and Data}
% ============================================================

[Place any additional figures, tables, or code listings that would interrupt the
main narrative but may be useful to a curious reader.]

\end{document}
